%!TEX root = ../main.tex
% =============================================================
%  CAPÍTULO 7 — PRESENTACIÓN DE RESULTADOS
%                                       (Entrega 7 · 5 % · 13/10)
% =============================================================

\chapter{Presentación de resultados}\label{cap:resultados}

\begin{guia}[Guía de la cátedra: qué se evalúa en este capítulo]
\begin{enumerate}
    \item \textbf{Requerimientos formales}: \gls{rf} y \gls{rnf} con ID,
          descripción, prioridad y criterio de aceptación medible,
          trazados a los indicadores del capítulo~\ref{cap:metodologia}.
    \item \textbf{Evidencia de implementación}: capturas reales del sistema
          funcionando, enlaces al repositorio (anexo~\ref{anx:repositorios})
          y al despliegue.
    \item \textbf{Resultados cuantitativos honestos}: diseño experimental,
          métricas, matrices de confusión o su equivalente en el dominio,
          incluyendo lo que salió mal y su análisis. Reportar también lo
          malo y argumentarlo vale más que esconderlo. Hiperparámetros y
          versiones en el apéndice~\ref{apx:tecnicos}.
    \item \textbf{Validación con usuarios y expertos}: idealmente
          triangulada (métricas técnicas, encuesta, entrevista) y
          traducida a decisiones de diseño concretas. Instrumentos en los
          apéndices \ref{apx:encuesta} y \ref{apx:entrevista}.
\end{enumerate}
El sistema tiene que existir y andar: la entrega final incluye el
proyecto funcional.
\end{guia}

\completar{Desarrollo del capítulo.}

\section{\completar{Nombre de la sección}}

\completar{Texto.}

\begin{table}[H]
    \centering
    \caption{Requerimientos funcionales (ejemplo de estructura)}
    \label{tab:rf}
    \small
    \begin{tabularx}{\textwidth}{l L L l L}
        \toprule
        \tb{ID} & \tb{Nombre} & \tb{Descripción} & \tb{Prioridad} & \tb{Criterio de aceptación} \\
        \midrule
        RF-01 & Autenticación & El sistema permite registrarse e iniciar sesión con credenciales propias & Alta & Un usuario nuevo completa el registro y accede en menos de 2 minutos \\
        RF-02 & \completar{...} & \completar{...} & \completar{...} & \completar{...} \\
        \bottomrule
    \end{tabularx}
\end{table}

\begin{table}[H]
    \centering
    \caption{Requerimientos no funcionales (ejemplo de estructura)}
    \label{tab:rnf}
    \small
    \begin{tabularx}{\textwidth}{l L L l L}
        \toprule
        \tb{ID} & \tb{Atributo} & \tb{Descripción} & \tb{Prioridad} & \tb{Umbral medible} \\
        \midrule
        RNF-01 & Rendimiento & Tiempo de respuesta de las consultas principales & Alta & Latencia p95 menor a 5 s (indicador I-02) \\
        RNF-02 & Calidad del modelo & Capacidad predictiva sobre datos no vistos & Alta & $F_1 \geq 0{,}70$ (indicador I-01) \\
        RNF-03 & \completar{...} & \completar{...} & \completar{...} & \completar{...} \\
        \bottomrule
    \end{tabularx}
\end{table}

\section{\completar{Nombre de la sección}}

\completar{Texto.}

\begin{figure}[H]
    \centering
    \includegraphics[width=0.85\textwidth]{example-image-c}
    \caption{Captura del sistema en funcionamiento (ejemplo de figura)}
    \label{fig:captura}
\end{figure}

\begin{table}[H]
    \centering
    \caption{Resultados sobre el conjunto de prueba (ejemplo de estructura)}
    \label{tab:metricas}
    \small
    \begin{tabular}{l r r r r}
        \toprule
        \tb{Modelo} & \tb{Precisión} & \tb{Exhaustividad} & \tb{$F_1$} & \tb{Tiempo (s)} \\
        \midrule
        Línea de base & \completar{...} & \completar{...} & \completar{...} & \completar{...} \\
        Modelo propuesto & \completar{...} & \completar{...} & \completar{...} & \completar{...} \\
        \bottomrule
    \end{tabular}
\end{table}

\begin{table}[H]
    \centering
    \caption{Traducción de la retroalimentación en decisiones de diseño (ejemplo de estructura)}
    \label{tab:feedback-decisiones}
    \small
    \begin{tabularx}{\textwidth}{l L L l}
        \toprule
        \tb{Fuente} & \tb{Hallazgo} & \tb{Decisión} & \tb{Estado} \\
        \midrule
        Encuesta (n=\completar{..}) & \completar{El 60\,\% no encontró la función X} & \completar{Mover X al menú principal} & Implementado \\
        Experto 1 & \completar{...} & \completar{...} & Planificado \\
        \bottomrule
    \end{tabularx}
\end{table}
